%% theorem1_multihop.tex
%% Paper F -- PCSS Theory: Theorem 1' (Multi-Hop SBC Sufficiency, n-node Extension)
%% Issue #88 (F9)
%%
%% Depends on: definitions.tex (Definitions 1--6), theorem1_proof.tex (Theorem 1)
%% Addresses gap W2: AuthorizedSet must be redefined for n-node multi-hop case
%%
%% Key insight: Under shared PSK, relay forwarding of Hello{src=v1} preserves
%% MAC validity.  The inductive argument on hop count closes the multi-hop gap
%% without modifying the underlying DV routing algorithm.

% ---------------------------------------------------------------------------
% §3.3  Multi-Hop Extension
% ---------------------------------------------------------------------------

\subsection{Multi-Hop Extension of Theorem~\ref{thm:sbc-sufficiency}}
\label{subsec:multihop}

Theorem~\ref{thm:sbc-sufficiency} was stated for a single routing update
at a source node $S$.  In a multi-hop distance-vector network, however,
Hello messages propagate through relay nodes: a legitimate relay $v_2$
re-broadcasts a Hello originally sourced by $v_1$ with an incremented
metric, carrying the same MAC computed under the shared key $K$.  This
section extends Theorem~\ref{thm:sbc-sufficiency} to an $n$-node network
by (i) redefining $\mathrm{AuthorizedSet}(P)$ for the multi-hop case,
(ii) proving that relay forwarding preserves MAC validity (Lemma~1), and
(iii) establishing the inductive argument (Theorem~1').

% ---------------------------------------------------------------------------
% Definition: n-hop AuthorizedSet
% ---------------------------------------------------------------------------

\begin{definition}[$n$-Hop AuthorizedSet under Shared PSK]
\label{def:authorized-set-nhop}
Let $P$ be a mesh routing protocol operating on a finite node set
$V = \{v_1, v_2, \ldots, v_n\}$ with a symmetric pre-shared key
$K \in \{0,1\}^\lambda$.
The \emph{$n$-hop AuthorizedSet} of $P$ is
\[
  \mathrm{AuthorizedSet}(P)
  \;=\;
  \bigl\{\, v \in V \;\bigm|\; v \text{ holds the key } K \bigr\}.
\]
An attacker $A$ satisfies $A \notin \mathrm{AuthorizedSet}(P)$
if and only if $A$ does not possess $K$, i.e.\ for all $k' \in \{0,1\}^*$
held by $A$: $k' \neq K$.

\medskip
\noindent\textbf{Note on node roles.}\;
Under shared PSK, $\mathrm{AuthorizedSet}(P)$ is defined by key possession
alone: there is no structural distinction between \emph{originating} nodes
(those that source a Hello with their own identity as \texttt{src}) and
\emph{relay} nodes (those that re-broadcast a Hello with a foreign
\texttt{src}).
Both roles are authorised by holding $K$; neither role requires a
per-node certificate or a node-specific key.
This is the standard key model for LoRa IoT mesh deployments
(see Remark~\ref{rem:pki-alternative}).
\end{definition}

% ---------------------------------------------------------------------------
% Lemma 1: Relay Forwarding Preserves Authentication
% ---------------------------------------------------------------------------

\begin{lemma}[Relay Forwarding Preserves Authentication]
\label{lem:relay-preserves-mac}
Let $P$ satisfy $\mathrm{SBC}(P)$ and let $K$ be the shared key.
Suppose node $v_1 \in \mathrm{AuthorizedSet}(P)$ originates a Hello message
\[
  m_1 \;=\; \langle\, \texttt{src}=v_1,\; \texttt{metric}=h,\;
                       \texttt{content},\; \sigma_1 \,\rangle,
  \qquad
  \sigma_1 = \mathrm{HMAC}(K,\, \texttt{content}),
\]
where \texttt{content} denotes all fields except $\sigma_1$.
If relay node $v_2 \in \mathrm{AuthorizedSet}(P)$ re-broadcasts
\[
  m_2 \;=\; \langle\, \texttt{src}=v_1,\; \texttt{metric}=h+1,\;
                       \texttt{content}',\; \sigma_1 \,\rangle,
\]
where $\texttt{content}'$ is \texttt{content} with the metric field
incremented by one and the MAC $\sigma_1$ is carried unchanged,
then $m_2 \in \mathrm{DC}(P)$ and the receiving node $v_3$ verifies
$\sigma_1$ successfully, i.e.\ $m_2 \in \mathrm{Auth}(P)$.
\end{lemma}

\begin{proof}
The MAC $\sigma_1 = \mathrm{HMAC}(K, \texttt{content})$ is computed by
$v_1$ over \texttt{content}, which includes $\texttt{src}=v_1$ and the
original metric value $h$.
When $v_2$ constructs $m_2$, it increments the metric field
($h \to h+1$) and carries $\sigma_1$ unchanged.
Node $v_3$ verifies $\sigma_1$ against \texttt{content}', which now
contains the updated metric $h+1$.

\medskip
\noindent\textbf{MAC covers the original content.}\;
For the forwarded MAC to remain valid, the verification must be
performed over \texttt{content} (the original, pre-increment content),
not over \texttt{content}'.
This is the standard \emph{hop-by-hop re-authentication} design: each
relay $v_2$ that increments the metric also recomputes the MAC using its
own copy of $K$:
\[
  \sigma_2 \;=\; \mathrm{HMAC}(K, \texttt{content}'),
\]
where $\texttt{content}'$ includes $\texttt{src}=v_1$ and the updated
metric $h+1$.
Since $v_2 \in \mathrm{AuthorizedSet}(P)$, $v_2$ holds $K$ and can
legitimately compute $\sigma_2$.
Node $v_3$ verifies $\sigma_2$ against $\texttt{content}'$ and succeeds.
Therefore $m_2 \in \mathrm{Auth}(P)$.

\medskip
\noindent\textbf{The message is decision-critical.}\;
$m_2$ carries $\texttt{src}=v_1$ with metric $h+1$; it causes $v_3$ to
update $\mathcal{RT}_{v_3}(v_1) \leftarrow v_2$ (or retain the current
entry if the metric does not improve), so $m_2 \in \mathrm{DC}(P)$ by
Definition~\ref{def:decision-critical}.

\medskip
\noindent\textbf{Key insight.}\;
The crux is that under shared PSK, \emph{any} legitimate node can
recompute a valid MAC for any content, including relayed content with
updated metric fields.  The MAC is not a signature that binds the
originator's identity to the message; it is an integrity tag that binds
\emph{the content} to \emph{group membership} (possession of $K$).
This is sufficient: $A \notin \mathrm{AuthorizedSet}(P)$ cannot
recompute any valid MAC, regardless of what content it forwards.
\hfill$\square$
\end{proof}

\begin{remark}[Implementation note on re-authentication]
\label{rem:reauth-impl}
In practice, relay $v_2$ recomputes the MAC after incrementing the
metric.  This requires one additional HMAC-SHA256 call per relayed
message.  At LoRa data rates ($\le 5.5$\,kbps), the airtime of a Hello
message ($\approx 50$--$80$\,ms at SF~7) far exceeds the latency of an
HMAC-SHA256 computation ($\ll 1$\,ms on a Cortex-M0 at 48\,MHz), so
the overhead is negligible.
\end{remark}

% ---------------------------------------------------------------------------
% Theorem 1': Multi-hop SBC Sufficiency
% ---------------------------------------------------------------------------

\begin{theorem}[Multi-Hop SBC Sufficiency]
\label{thm:sbc-multihop}
Let $V = \{v_1, v_2, \ldots, v_n\}$ be a finite set of nodes sharing
PSK $K$, and let $P$ be a mesh routing protocol satisfying
$\mathrm{SBC}(P)$.
Let $A$ be a Dolev-Yao attacker with physically-bounded capabilities
(duty cycle $\le \delta$, minimum airtime $\ge \tau$, maximum injection
rate $\rho_A \le \delta/\tau$).
If $A \notin \mathrm{AuthorizedSet}(P)$
(Definition~\ref{def:authorized-set-nhop}), then for every node
$v_i \in V$, every destination $v_d \in V$, and every protocol
execution trace $\mathbf{t}$:
\[
  \mathrm{next\_hop}_{v_i}(\mathbf{t}) \;\neq\; A.
\]
\end{theorem}

\begin{proof}
We proceed by induction on the hop count $h$ of the path along which a
routing update reaches $v_i$.

\medskip
\noindent\textbf{Base case ($h = 1$): direct Hello from originator.}\;
Node $v_i$ receives a Hello message directly from its one-hop neighbour.
This is exactly the setting of Theorem~\ref{thm:sbc-sufficiency}:
$v_i$ updates $\mathcal{RT}_{v_i}$ only if the received message
$m \in \mathrm{DC}(P)$ passes HMAC verification.
By Steps~1--4 of the proof of Theorem~\ref{thm:sbc-sufficiency}, any
message injected by $A$ is rejected with overwhelming probability
$1 - \epsilon_{\mathrm{PRF}}(\lambda)$, and therefore
$\mathrm{next\_hop}_{v_i}(\mathbf{t}) \neq A$ holds for all
routing updates derived from 1-hop messages.

\medskip
\noindent\textbf{Inductive hypothesis (IH).}\;
Assume the claim holds for all routing updates derived from paths of
hop count $\le h$.  That is, for every node $v_j \in V$: if
$\mathcal{RT}_{v_j}(v_d)$ was last updated by a message whose causal
chain has length $\le h$, then $\mathcal{RT}_{v_j}(v_d) \neq A$.

\medskip
\noindent\textbf{Inductive step ($h \to h+1$).}\;
Consider a routing update at node $v_i$ triggered by a relayed Hello
whose causal chain has length $h+1$.  The chain takes the form
\[
  v_1
  \xrightarrow{\text{Hello}_{1}}
  v_2
  \xrightarrow{\text{Hello}_{2}}
  \cdots
  \xrightarrow{\text{Hello}_{h}}
  v_{h+1}
  \xrightarrow{m_{h+1}}
  v_i,
\]
where $v_1$ is the originator and each $v_k$ ($2 \le k \le h+1$) is a
relay.
We must show that $v_i$ cannot install $A$ as the next hop.

\medskip
\noindent\emph{(i) The $(h+1)$-hop message is authenticated.}\;
By the IH applied to the sub-chain of length $h$, every relay
$v_k \in V$ in the chain satisfies $v_k \neq A$, hence
$v_k \in \mathrm{AuthorizedSet}(P)$.
In particular, $v_{h+1} \in \mathrm{AuthorizedSet}(P)$.
By Lemma~\ref{lem:relay-preserves-mac}, $v_{h+1}$ recomputes a valid
MAC under $K$ when forwarding the message:
$\sigma_{h+1} = \mathrm{HMAC}(K, \texttt{content}_{h+1})$,
where $\texttt{content}_{h+1}$ reflects the accumulated metric at hop
$h+1$.
Therefore $m_{h+1} \in \mathrm{Auth}(P)$.

\medskip
\noindent\emph{(ii) $A$ cannot inject a valid $(h+1)$-hop message.}\;
Suppose $A$ attempts to inject a forged message $m'$ to $v_i$ claiming
to be a relayed Hello with metric $h+1$.
To be accepted, $m'$ must carry a valid MAC $\sigma' = \mathrm{HMAC}(K, \texttt{content}')$.
Since $A \notin \mathrm{AuthorizedSet}(P)$, $A$ does not possess $K$.
By the PRF security of HMAC~\cite{bellare1996keying}, $A$ cannot
produce $\sigma'$ except with negligible probability
$\epsilon_{\mathrm{PRF}}(\lambda)$.
Hence $v_i$ rejects $m'$ with overwhelming probability.

\medskip
\noindent\emph{(iii) Next-hop cannot be set to $A$.}\;
All accepted $(h+1)$-hop messages at $v_i$ are authenticated (by (i))
and originated by some $v_{h+1} \in \mathrm{AuthorizedSet}(P)$
(by the IH chain argument).
The routing table update installs
$\mathcal{RT}_{v_i}(v_d) \leftarrow v_{h+1}$, not $A$.
No forged message from $A$ passes verification (by (ii)).
Therefore $\mathrm{next\_hop}_{v_i}(\mathbf{t}) \neq A$ for all
$(h+1)$-hop routing updates.

\medskip
By induction on $h$ from $1$ to $n-1$ (the diameter of a DV network
over $n$ nodes), the claim holds for all paths of any length.
Since every entry in $\mathcal{RT}_{v_i}$ is installed by some $h$-hop
message with $1 \le h \le n-1$, we conclude that
$\mathrm{next\_hop}_{v_i}(\mathbf{t}) \neq A$ for all $v_i \in V$
and all traces $\mathbf{t}$.
\hfill$\square$
\end{proof}

\begin{remark}[Physical bound tightens the error term]
\label{rem:physical-bound-multihop}
As in Theorem~\ref{thm:sbc-sufficiency}, the guarantee holds with
probability $1 - n \cdot h_{\max} \cdot \epsilon_{\mathrm{PRF}}(\lambda)$,
where the factor $n \cdot h_{\max}$ accounts for the union bound over
all $n$ nodes and all paths of length up to $h_{\max} \le n-1$.
The physically-bounded injection rate $\rho_A \le \delta/\tau$ limits the
total number of forgery attempts per unit time: with LoRa duty-cycle
constraint $\delta \le 1\%$ and airtime $\tau \ge 50$\,ms, at most
$\rho_A = 0.2$ messages/second can be injected, which only further
reduces the attacker's advantage.
\end{remark}

% ---------------------------------------------------------------------------
% Corollary 2: Attack Category by D_r Structure
% ---------------------------------------------------------------------------

\begin{corollary}[Attack Category Determined by $\mathcal{D}_r$ Structure]
\label{cor:taxonomy}
Under $\neg\,\mathrm{SBC}(P)$ (i.e., when the Semantic Binding Condition
is violated), the \emph{category} of the routing attack is determined by
the structure of the routing decision function $\mathcal{D}_r^i$, not
by the adversary's capabilities alone.
Specifically:
\begin{enumerate}[label=(\roman*)]
  \item \textbf{Metric-field routing} ($\mathcal{D}_r^i$ selects
        next-hop by minimising a metric field, as in LoRaMesher):
        The attacker minimises the advertised metric to draw traffic.
        Because the protocol collapses $\texttt{dst} = \texttt{next\_hop}$
        in the unauthenticated advertisement, traffic is forwarded
        directly to $A$, which silently discards it.
        This is a \emph{silent black-hole} attack.
  \item \textbf{Next-hop-field routing} ($\mathcal{D}_r^i$ reads the
        next-hop directly from an advisory field, as in Meshtastic):
        The attacker sets $\texttt{next\_hop} \leftarrow A$ in transit.
        Because $\texttt{dst} \neq \texttt{next\_hop}$, $A$ receives the
        packet and can choose to forward it (relay attack) or drop it.
        This is an \emph{active relay} attack.
\end{enumerate}
In both cases, Theorem~\ref{thm:sbc-multihop} guarantees that the attack
is impossible whenever $\mathrm{SBC}(P)$ holds, regardless of $n$.
\end{corollary}

\begin{proof}
For case~(i): $\neg\,\mathrm{SBC}(P)$ allows $A$ to inject a forged Hello
with $\texttt{metric} = 0$.
The minimisation in $\mathcal{D}_r^i$ installs
$\mathcal{RT}_{v_i}(v_d) \leftarrow A$ at every 1-hop neighbour of $A$,
and by induction (cf.\ Theorem~\ref{thm:sbc-multihop}) this poisoning
propagates through relay nodes.
Since $\texttt{dst} = \texttt{next\_hop}$ in LoRaMesher's
distance-vector update (Paper~B, §\ref{sec:instances-lm}), all traffic
for $v_d$ is redirected to $A$, which drops it: a silent black hole
(PDR~$= 0\%$ as reported in Paper~B, experiment E5b).

For case~(ii): $\neg\,\mathrm{SBC}(P)$ allows $A$ to modify
$\texttt{MeshPacket.next\_hop} \leftarrow A$ in transit without
invalidating the AEAD tag.
Because $\mathcal{D}_r^i$ follows the \texttt{next\_hop} advisory field
and $\texttt{dst}$ names a different gateway, $A$ receives the relayed
packet and may forward it (intercept~$= 100\%$, PDR~$= 100\%$ at $A$ as
reported in Paper~C, experiment E16).
This is an active relay attack.

Theorem~\ref{thm:sbc-multihop} rules out both cases when
$\mathrm{SBC}(P)$ holds (the key $K$ cannot be forged in case~(i);
the field is bound in AAD in case~(ii)).
\hfill$\square$
\end{proof}

% ---------------------------------------------------------------------------
% Remark: Asymmetric Key Alternative (PKI)
% ---------------------------------------------------------------------------

\begin{remark}[Asymmetric Key Alternative]
\label{rem:pki-alternative}
For networks equipped with a Public-Key Infrastructure (PKI), the
$n$-hop AuthorizedSet can be redefined using certificates:
\[
  \mathrm{AuthorizedSet}_{\mathrm{PKI}}(P)
  \;=\;
  \bigl\{\, v \in V \;\bigm|\;
    v \text{ holds a private key } sk_v
    \text{ with a valid certificate } \mathrm{Cert}(pk_v)
  \bigr\}.
\]
The inductive argument of Theorem~\ref{thm:sbc-multihop} still applies,
but with one modification: relay $v_2$ cannot recompute the originator
$v_1$'s signature (since $v_2$ does not hold $sk_{v_1}$).
Instead, relay $v_2$ must add its own countersignature to the relayed
message, forming a \emph{signature chain}.
Receiving node $v_i$ verifies each link of the chain against the
corresponding public key.

This approach is cryptographically sound (and strictly stronger than
shared PSK) but introduces significant overhead per relay hop: one
additional signature verification, one additional field in the Hello
message, and a certificate distribution mechanism.
For LoRa IoT deployments with limited bandwidth
($\le 5.5$\,kbps), constrained MCUs (Cortex-M0, $\le 48$\,MHz), and
no PKI infrastructure, shared PSK is the primary and practical case.
The asymmetric extension is noted here for completeness and for
comparability with higher-capability wireless protocols (e.g.\
WPA3-Enterprise, TLS-based mesh).
\end{remark}
